\documentclass[11pt,letterpaper]{article}
\usepackage[margin=1in]{geometry}
\usepackage{booktabs}
\usepackage{array}
\usepackage{enumitem}
\usepackage{xcolor}
\usepackage{colortbl}
\usepackage{graphicx}
\usepackage{fancyhdr}
\usepackage{titlesec}
\usepackage{hyperref}
\usepackage{tcolorbox}
\usepackage{multirow}
\usepackage{tabularx}

\hypersetup{colorlinks=true, linkcolor=blue!60!black, urlcolor=blue!60!black}

\pagestyle{fancy}
\fancyhf{}
\rhead{BME Lab 3 -- DIY}
\lhead{Music \& Involuntary Muscle Tension}
\rfoot{Page \thepage}

\titleformat{\section}{\large\bfseries}{}{0em}{}[\titlerule]
\titleformat{\subsection}{\normalsize\bfseries}{}{0em}{}

\newtcolorbox{warningbox}{colback=red!5!white, colframe=red!60!black,
  title={\textbf{Safety Warning}}, fonttitle=\bfseries}
\newtcolorbox{infobox}{colback=blue!5!white, colframe=blue!50!black,
  title={\textbf{Note}}, fonttitle=\bfseries}
\newtcolorbox{tipbox}{colback=green!5!white, colframe=green!40!black,
  title={\textbf{Tip}}, fonttitle=\bfseries}

\begin{document}

% ----------------------------------------------------------------
\begin{center}
{\LARGE \textbf{Laboratory 3: Do It Yourself}} \\[6pt]
{\Large Effect of Music Genre on Involuntary Muscle Tension (EMG)} \\[10pt]
{\normalsize BME -- Biosignal Acquisition \& Analysis} \\[4pt]
{\normalsize Protocol \& Procedure Guide}
\end{center}
\vspace{0.5cm}

% ----------------------------------------------------------------
\section{Hypothesis}

\textit{If a subject listens to high-arousal music (heavy metal) while at rest, then the RMS amplitude of their surface EMG---particularly in the upper trapezius (a postural ``stress muscle'')---will increase relative to a silent baseline, indicating involuntary muscle tension driven by sympathetic arousal.  Low-arousal music (classical) will produce equal or lower tension than baseline due to parasympathetic-mediated relaxation.  Therefore, music genre modulates resting skeletal-muscle tone through the autonomic nervous system.}

% ----------------------------------------------------------------
\section{Background \& Rationale}

Surface electromyography (EMG) measures the sum of motor unit action potentials (MUAPs) from underlying muscle fibres.  Even at ``rest,'' skeletal muscles maintain a baseline level of tonic activity driven by the gamma motor neuron loop.  This resting tone is modulated by descending input from the brainstem reticular formation, which itself is influenced by arousal level.

\textbf{Why music affects muscle tension:}
\begin{itemize}[nosep]
    \item Music activates the limbic system (amygdala, hypothalamus), which regulates the autonomic nervous system.
    \item Increased sympathetic arousal raises reticulospinal drive, elevating baseline motor neuron excitability and therefore resting muscle tone.
    \item The \textbf{upper trapezius} is particularly sensitive to psychological stress---it is the classic ``stress muscle'' that people tense involuntarily when anxious or aroused.
    \item The \textbf{forearm flexors} serve as a comparison: a distal limb muscle less directly coupled to stress-related postural tone.
\end{itemize}

\textbf{Why two muscles?}  By recording both trapezius and forearm, we can test whether the effect is global (both muscles) or preferentially affects postural/stress-related musculature.

\textbf{Why these three genres?}  Same arousal--valence logic as any music psychophysiology study:
\begin{itemize}[nosep]
    \item \textbf{Classical (Debussy):} Low arousal, positive valence --- expected to relax
    \item \textbf{Heavy metal (Metallica):} High arousal, mixed valence --- expected to increase tension
    \item \textbf{Pop/funk (Earth, Wind \& Fire):} Moderate--high arousal, positive valence --- moderate effect
\end{itemize}

% ----------------------------------------------------------------
\section{Equipment Checklist}

\begin{center}
\begin{tabular}{@{} c l l @{}}
\toprule
$\square$ & BioRadio device & charged, turned ON \\
$\square$ & 5 snap leads & 2 for CH1, 2 for CH2, 1 for GND \\
$\square$ & 5 snap electrodes (adhesive) & gel-type recommended \\
$\square$ & Alcohol wipes & for skin prep \\
$\square$ & Lab laptop with BioCapture & Design Studio computer \\
$\square$ & AirPods Pro (or ANC over-ear headphones) & see Section~\ref{sec:audio} \\
$\square$ & iPhone / phone with dB meter & for volume calibration \\
$\square$ & Stopwatch / timer & phone timer or watch \\
$\square$ & Pen + printed copy of this protocol & for timestamp log \\
$\square$ & USB drive or cloud & to transfer exported CSV after \\
\bottomrule
\end{tabular}
\end{center}

% ----------------------------------------------------------------
\section{Audio Setup}
\label{sec:audio}

\subsection{Headphone Selection (in order of preference)}
\begin{enumerate}[nosep]
    \item \textbf{AirPods Pro:} Active Noise Cancellation \textbf{ON}, Transparency \textbf{OFF}.
    \item \textbf{Over-ear closed-back ANC headphones} (Sony, Bose, etc.): ANC \textbf{ON}.
    \item \textbf{Over-ear closed-back passive headphones}: acceptable if the lab is quiet.
\end{enumerate}

\begin{warningbox}
\textbf{Do NOT use:} open-back headphones, earbuds without ANC, laptop speakers, or Transparency mode.  Ambient noise is an uncontrolled stimulus that will confound your results.
\end{warningbox}

\subsection{Volume Calibration --- 70\,dB SPL}

\begin{tcolorbox}[colback=yellow!5!white, colframe=orange!60!black, title=\textbf{How to set 70\,dB with AirPods + iPhone}]
\begin{enumerate}[nosep]
    \item Open \textbf{Settings $\rightarrow$ Control Center}.  Add \textbf{Hearing} if not already there.
    \item Connect AirPods.  Play 30\,s of any experimental track.
    \item Swipe into Control Center, tap the \textbf{ear icon} (Hearing).
    \item Adjust volume until the reading is \textbf{70\,dB}.
    \item \textbf{Lock the volume.}  Do not touch it for the entire experiment.
\end{enumerate}
\end{tcolorbox}

\begin{infobox}
\textbf{No AirPods / no iOS?} Use the free \textbf{NIOSH SLM} app (iOS) or \textbf{Sound Meter} (Android).  Hold the phone mic at the headphone driver, play a test track, adjust to 70\,dB.
\end{infobox}

\textbf{Same dB across all genres.} Do not adjust volume between songs. 70\,dB is a comfortable level that avoids startle confounds (which would cause EMG artifacts).

% ----------------------------------------------------------------
\section{Song Selection}

\begin{center}
\renewcommand{\arraystretch}{1.3}
\begin{tabularx}{\textwidth}{@{} l l X l @{}}
\toprule
\textbf{Condition} & \textbf{Song} & \textbf{Why} & \textbf{Start at} \\
\midrule
Calm & Debussy -- \textit{Clair de Lune} & Slow tempo ($\sim$72\,BPM), instrumental, consonant. & 0:00 \\
Intense & Metallica -- \textit{Master of Puppets} & Fast tempo ($\sim$212\,BPM), distorted, aggressive. & 0:00 \\
Happy & Earth, Wind \& Fire -- \textit{September} & Major key, danceable ($\sim$126\,BPM), positive. & 0:00 \\
\bottomrule
\end{tabularx}
\end{center}

\begin{tipbox}
Pre-download all three songs for offline playback (no ads/buffering).  Queue them in order: Debussy $\rightarrow$ Metallica $\rightarrow$ September.  Put your phone in \textbf{Do Not Disturb}.
\end{tipbox}

% ----------------------------------------------------------------
\section{BioRadio \& Electrode Setup}

This follows the same general procedure as Lab~2 Path~2, with different electrode placement.

\subsection{BioCapture Configuration}

\begin{enumerate}[nosep]
    \item Open BioCapture.  \textbf{Do not} connect the BioRadio to the computer yet.
    \item Turn the BioRadio ON.
    \item In BioCapture: \textbf{Device $\rightarrow$ Select Device $\rightarrow$} your group's BioRadio $\rightarrow$ \textbf{Connect}.
    \item Click \textbf{Device Config}:
    \begin{itemize}[nosep]
        \item Check \textbf{CH1}.  Name it \textbf{``Forearm''}.  Type: \textbf{EMG}.
        \item Check \textbf{CH2}.  Name it \textbf{``Trapezius''}.  Type: \textbf{EMG}.
        \item Uncheck any other channels.
        \item Set sampling rate to \textbf{2\,kHz}.
    \end{itemize}
    \item Click \textbf{Program Device}.
\end{enumerate}

\subsection{Electrode Placement}

You need \textbf{5 electrodes}: 2 for CH1, 2 for CH2, 1 ground.

\begin{center}
\renewcommand{\arraystretch}{1.3}
\begin{tabular}{@{} l l p{7cm} @{}}
\toprule
\textbf{Channel} & \textbf{Muscle} & \textbf{Electrode Placement} \\
\midrule
CH1 (+/--) & Forearm flexors & Two electrodes on the \textbf{center of the inner forearm}, $\sim$3\,cm apart, along the muscle belly.  Same placement as Lab~2. \\
\addlinespace
CH2 (+/--) & Upper trapezius & Two electrodes on the \textbf{upper trapezius}, midway between the C7 vertebra (the bony bump at the base of your neck) and the acromion (tip of your shoulder).  Place them $\sim$3\,cm apart, along the muscle fibres (i.e., along the line from neck to shoulder). \\
\addlinespace
GND & --- & One electrode on the \textbf{bony part of the elbow} (olecranon). \\
\bottomrule
\end{tabular}
\end{center}

\textbf{Skin preparation:}
\begin{itemize}[nosep]
    \item Clean each electrode site with an alcohol wipe.  Let it dry for 10\,seconds.
    \item This removes oils and dead skin, improving electrode--skin contact and reducing impedance.
\end{itemize}

\textbf{Wire connections (same colour scheme as Lab~2):}
\begin{itemize}[nosep]
    \item CH1 +/-- (black and red plugs next to blue GND) $\rightarrow$ forearm electrodes
    \item CH2 +/-- (red and black plugs next to CH1) $\rightarrow$ trapezius electrodes
    \item GND (blue plug) $\rightarrow$ elbow electrode
\end{itemize}

\subsection{Signal Check}

\begin{enumerate}[nosep]
    \item Click \textbf{Run} in BioCapture.  Set Graph Timespan to \textbf{10\,s}.
    \item Right-click the y-axis of each graph $\rightarrow$ \textbf{Y-Axis $\rightarrow$ Auto Scale Quick}.
    \item You should see a low-amplitude noisy baseline on both channels.
    \item \textbf{Test:} briefly shrug your shoulders---CH2 (trapezius) should spike.  Make a fist---CH1 (forearm) should spike.
    \item If a channel is flat or railed, check the corresponding electrode connections and skin prep.
\end{enumerate}

% ----------------------------------------------------------------
\section{Experimental Protocol}
\label{sec:protocol}

\textbf{Total duration: $\sim$20 minutes of recording + 5 minutes setup.}

\subsection{Pre-experiment Checklist}

\begin{enumerate}[nosep]
    \item[$\square$] BioRadio connected, channels configured, signal verified (Section~6).
    \item[$\square$] Songs downloaded and queued on phone.
    \item[$\square$] AirPods connected, ANC ON, volume set to 70\,dB.
    \item[$\square$] Stopwatch ready.
    \item[$\square$] Printed protocol + pen in front of you.
    \item[$\square$] Subject seated comfortably: feet flat on floor, forearms resting on table, shoulders relaxed (not hunched).
\end{enumerate}

\subsection{Recording Procedure}

Your lab partner operates the timer, music controls, and BioCapture recording.  The subject (you) sits still.

\begin{enumerate}
    \item \textbf{Start recording:} In BioCapture, click the \textbf{Record} button (red dot).  Save the file as \texttt{music\_emg}. Recording begins automatically.
    \item \textbf{Start the stopwatch} at the same moment.  Write down the start time below.
\end{enumerate}

\begin{center}
\renewcommand{\arraystretch}{1.4}
\begin{tabular}{@{} >{\bfseries}c l c c p{5.5cm} @{}}
\toprule
\textbf{\#} & \textbf{Condition} & \textbf{Duration} & \textbf{Stopwatch} & \textbf{Instructions} \\
\midrule
\rowcolor{gray!10}
1 & Silence (Baseline) & 3 min & 0:00 & Sit still, eyes open, breathe normally.  No talking, no phone. \\
2 & Silence (Recovery) & 2 min & 3:00 & Continue sitting still. \\
\rowcolor{blue!5}
3 & Calm Music & 3 min & 5:00 & Partner plays \textit{Clair de Lune}. Listen passively. \\
4 & Silence (Recovery) & 2 min & 8:00 & Partner pauses music. \\
\rowcolor{red!5}
5 & Intense Music & 3 min & 10:00 & Partner plays \textit{Master of Puppets}. Listen passively. \\
6 & Silence (Recovery) & 2 min & 13:00 & Partner pauses music. \\
\rowcolor{yellow!10}
7 & Happy Music & 3 min & 15:00 & Partner plays \textit{September}. Listen passively. \\
8 & Silence (Recovery) & 2 min & 18:00 & Partner pauses music. Final cooldown. \\
\midrule
& \textbf{STOP} & & 20:00 & Partner clicks \textbf{Record} again to stop.  Data auto-saves. \\
\bottomrule
\end{tabular}
\end{center}

\vspace{0.3cm}
\textbf{Timestamp log} (write actual stopwatch times here in case of delays):

\begin{center}
\renewcommand{\arraystretch}{1.5}
\begin{tabular}{@{} l c @{}}
\toprule
\textbf{Event} & \textbf{Actual stopwatch time} \\
\midrule
Recording started (Record button pressed)                   & \underline{\hspace{3cm}} \\
Baseline started (Phase 1)                                  & \underline{\hspace{3cm}} \\
Recovery 1 started (Phase 2)                                & \underline{\hspace{3cm}} \\
Calm music started -- \textit{Clair de Lune} (Phase 3)      & \underline{\hspace{3cm}} \\
Recovery 2 started (Phase 4)                                & \underline{\hspace{3cm}} \\
Intense music started -- \textit{Master of Puppets} (Phase 5) & \underline{\hspace{3cm}} \\
Recovery 3 started (Phase 6)                                & \underline{\hspace{3cm}} \\
Happy music started -- \textit{September} (Phase 7)         & \underline{\hspace{3cm}} \\
Recovery 4 started (Phase 8)                                & \underline{\hspace{3cm}} \\
Recording stopped                                           & \underline{\hspace{3cm}} \\
\bottomrule
\end{tabular}
\end{center}

\subsection{Critical Rules During Recording}

\begin{itemize}[nosep]
    \item \textbf{Do not move.}  Any movement creates massive EMG artefacts that dwarf the resting-tension signal.  Keep arms flat on the table, shoulders relaxed.
    \item \textbf{Do not talk, swallow frequently, or clench your jaw.}  These activate neck/facial muscles that can couple into the trapezius signal.
    \item \textbf{Do not look at the BioCapture screen.}  Seeing your own signals can cause anxiety (a confound).  Have the screen face your partner.
    \item \textbf{Breathe normally.}  Deep breathing changes thoracic muscle activity.
    \item \textbf{Do not tap, bob your head, or move to the music.}
\end{itemize}

% ----------------------------------------------------------------
\section{Data Export}

After the experiment:

\begin{enumerate}[nosep]
    \item In BioCapture: \textbf{Open $\rightarrow$ Recording\ldots} $\rightarrow$ select \texttt{music\_emg}.
    \item Click \textbf{Data $\rightarrow$ Export to .CSV File}.  Name it \texttt{lab3\_emg\_data.csv}.  Click Export.
    \item Transfer the CSV to the computer where you will run the analysis (USB drive, email, cloud, etc.).
    \item Place the CSV in the same folder as \texttt{analyze\_emg.py}.
\end{enumerate}

% ----------------------------------------------------------------
\section{Data Analysis}

\subsection{Setup (one-time)}
\begin{enumerate}[nosep]
    \item Install Python dependencies: \texttt{pip3 install numpy scipy matplotlib pandas}
    \item Open \texttt{analyze\_emg.py} in a text editor.
    \item Update \texttt{PHASE\_START\_S} with your actual stopwatch timestamps (convert MM:SS to seconds).
    \item Verify the column indices match your CSV.  Biocapture typically exports: column~0 = Time, column~1 = CH1, column~2 = CH2.  If your CSV differs, update \texttt{TIME\_COL}, \texttt{CH1\_COL}, \texttt{CH2\_COL}.
    \item Run: \texttt{python3 analyze\_emg.py}
\end{enumerate}

\subsection{What the Script Produces}

\begin{enumerate}[nosep]
    \item \textbf{fig1\_emg\_timeseries.png} --- RMS envelope of both channels over the full 20 minutes, with phase markers.  This is the main result figure.
    \item \textbf{fig2\_emg\_comparison.png} --- Bar charts of mean RMS (mV) per condition for forearm and trapezius, with standard deviation error bars.
    \item \textbf{fig3\_emg\_snippets.png} --- 2-second raw EMG snippets from each condition, showing the waveform character (amplitude, burst activity).
    \item \textbf{fig4\_emg\_spectrum.png} --- Power spectral density of the EMG for each condition.  Useful for verifying signal quality (should see power in the 20--450\,Hz band, not dominated by 60\,Hz noise).
    \item \textbf{Console table} --- Mean $\pm$ SD of RMS amplitude per phase.
\end{enumerate}

\subsection{Signal Processing Details (for your report)}

\begin{itemize}[nosep]
    \item \textbf{Bandpass filter:}  4th-order Butterworth, 20--450\,Hz.  The 20\,Hz high-pass removes motion artefact and DC offset.  The 450\,Hz low-pass is below the Nyquist frequency ($F_s/2 = 1000$\,Hz).
    \item \textbf{RMS envelope:}  Root-mean-square computed in a 250\,ms sliding window.  RMS is the standard measure of EMG amplitude because it is proportional to the number of active motor units and their firing rates.
    \item \textbf{Power spectrum:}  FFT of each phase's filtered EMG.  The median frequency of the power spectrum can indicate fatigue (shifts lower), though this is less relevant for resting measurements.
\end{itemize}

% ----------------------------------------------------------------
\section{Expected Results}

\begin{center}
\renewcommand{\arraystretch}{1.3}
\begin{tabular}{@{} l c c @{}}
\toprule
\textbf{Condition vs.\ Baseline} & \textbf{Forearm EMG RMS} & \textbf{Trapezius EMG RMS} \\
\midrule
Calm Music     & no change or slight decrease & decrease (relaxation) \\
Intense Music  & slight increase & \textbf{increase} (stress response) \\
Happy Music    & slight increase & slight increase \\
\bottomrule
\end{tabular}
\end{center}

The trapezius is expected to show the largest effect because it is a postural muscle strongly coupled to stress/arousal pathways (reticulospinal tract).  The forearm, being a distal voluntary-use muscle, may show a smaller or no effect.

If your results differ, discuss possible reasons in your reflection: individual differences in music preference, electrode placement quality, difficulty staying perfectly still, etc.

% ----------------------------------------------------------------
\section{Report Structure}

Submit a single PDF containing:

\begin{enumerate}[nosep]
    \item \textbf{Background \& Rationale} ($\sim$0.5 page): How music affects the nervous system and why it might modulate resting muscle tension.
    \item \textbf{Hypothesis:} The if/then/therefore statement.
    \item \textbf{Methods} ($\sim$0.5 page): BioRadio setup, electrode placement, audio setup (headphone type, dB), song choices, protocol timeline.
    \item \textbf{Results} ($\sim$1 page): Figures 1--4 with descriptions.
    \item \textbf{Reflection} ($\sim$0.5 page): Did results support the hypothesis?  What would you change?  Limitations (single subject, no counterbalancing, genre is a compound variable, etc.).
\end{enumerate}

% ----------------------------------------------------------------
\section{Troubleshooting}

\begin{tabularx}{\textwidth}{@{} l X @{}}
\toprule
\textbf{Problem} & \textbf{Solution} \\
\midrule
BioRadio won't connect & Turn off and on.  Make sure Bluetooth is enabled on the laptop.  Try Device $\rightarrow$ Select Device again. \\
\addlinespace
Signal is flat on one channel & Check snap lead is firmly connected to the electrode.  Check the electrode is making good skin contact (press down edges).  Re-prep skin with alcohol wipe. \\
\addlinespace
Large 60\,Hz oscillation visible & Move laptop charger away from electrodes.  Ensure ground electrode has good contact.  The bandpass filter in the analysis script will attenuate this, but better to minimize at source. \\
\addlinespace
Signal clips / rails to max & Right-click y-axis $\rightarrow$ Auto Scale Quick.  If it persists, the electrode may be poorly placed or the muscle is being contracted. \\
\addlinespace
CSV has unexpected columns & Open the CSV in a text editor.  Update the column indices (\texttt{TIME\_COL}, \texttt{CH1\_COL}, \texttt{CH2\_COL}) in \texttt{analyze\_emg.py} to match. \\
\addlinespace
Music interrupted by notification & Put phone in Do Not Disturb / Focus mode before starting. \\
\bottomrule
\end{tabularx}

\vfill
\begin{center}
\small\textit{Protocol prepared for BME Lab 3 -- DIY.}
\end{center}

\end{document}
